% =====================================================================
% CV - Jiajun Hu
% Source file: files/CV-JiajunHu-EN.tex   (Chinese version: CV-JiajunHu-CH.tex)
% Compile: latexmk -xelatex CV-JiajunHu-EN.tex   (or: latexmk from files/)
% Dependencies: standard TeX Live (fontspec, xcolor, microtype)
% =====================================================================
% !TEX program = xelatex
\documentclass[10pt,letterpaper]{article}

\usepackage[margin=0.6in,top=0.55in,bottom=0.55in]{geometry}
\usepackage{fontspec}           % XeLaTeX font handling
\setmainfont{Charter}           % macOS system serif
\usepackage{microtype}          % better justification and spacing
\usepackage{xcolor}
\usepackage{hyperref}

% --- Palette -----------------------------------------------------------
\definecolor{navy}{RGB}{22,45,95}
\definecolor{gray}{RGB}{95,95,95}

\hypersetup{colorlinks=true,urlcolor=navy,linkcolor=black,pdftitle={Jiajun Hu CV}}

\pagestyle{empty}
\raggedbottom                 % do not stretch glue to fill the page
\setlength{\parindent}{0pt}
\setlength{\parskip}{0pt}       % vertical spacing is set explicitly by the macros below
\setlength{\emergencystretch}{1.5em}  % lets long venue names and unbreakable links wrap

% --- Type scale: section > entry title > detail lines and bullets -------
\newcommand{\bodysize}{\fontsize{10}{13}\selectfont}      % summary text
\newcommand{\titlesize}{\fontsize{10.5}{13}\selectfont}   % school, employer, project titles
\newcommand{\detailsize}{\fontsize{9.5}{12}\selectfont}   % degrees, venues, dates, bullets
\AtBeginDocument{\bodysize}

% --- Section heading --------------------------------------------
\newcommand{\cvsection}[1]{%
  \needspace{6\baselineskip}\vspace{9pt}\nopagebreak
  {\fontsize{12.5}{15}\selectfont\bfseries\color{navy}#1\par}\vspace{2pt}\nopagebreak
  {\color{navy}\hrule width\linewidth height0.7pt}\nopagebreak
  \vspace{4pt}\nopagebreak}

% --- Entry: bold title left, gray location right; \subentry: smaller detail line below ---
\newcommand{\entry}[2]{\par\vspace{3pt}{\titlesize\textbf{#1}\hfill{\detailsize\color{gray}#2}\par}\nopagebreak}
\newcommand{\subentry}[2]{{\detailsize\emph{#1}\hfill{\color{gray}#2}\par}\nopagebreak}

% --- Break the page early if less than #1 of space remains (minimal needspace) ---
\makeatletter
\newcommand{\needspace}[1]{\par\penalty-100\begingroup
  \setlength{\dimen@}{#1}\dimen@ii\pagegoal\advance\dimen@ii-\pagetotal
  \ifdim\dimen@>\dimen@ii\ifdim\dimen@ii>\z@\vfil\fi\penalty-\@M\fi\endgroup}
\makeatother

% --- Project inside an entry: bold title, venue in gray italics below ---
\newcommand{\project}[2]{\needspace{7\baselineskip}\vspace{7pt}{\titlesize\textbf{#1}\par}\nopagebreak{\detailsize\color{gray}\emph{#2}\par}\nopagebreak}

% --- Bullet list: smaller than titles, with air between items ------------
\newenvironment{cvlist}{%
  \detailsize
  \begin{list}{\textbullet}{%
    \setlength{\itemsep}{4.5pt}\setlength{\parsep}{0pt}%
    \setlength{\topsep}{3pt}\setlength{\partopsep}{0pt}%
    \setlength{\leftmargin}{1.2em}\setlength{\labelwidth}{0.7em}%
    \setlength{\labelsep}{0.5em}}}%
  {\end{list}}

\begin{document}

% =====================================================================
% HEADER
% =====================================================================
{\centering
  {\fontsize{22}{26}\selectfont\bfseries\color{navy}Jiajun Hu}\\[6pt]
  \detailsize\color{gray}(412) 726-8574\ \ \textbar\ \
  \href{mailto:jiajunh5@asu.edu}{jiajunh5@asu.edu}\ \ \textbar\ \
  \href{https://codingmmax.github.io/}{codingmmax.github.io}\par}
\vspace{6pt}

% =====================================================================
% SUMMARY
% =====================================================================
\cvsection{Summary}

Ph.D. student in Computer Engineering working on FPGA architecture \& CAD, in-memory computing, and \mbox{DRAM-PIM} for ML/LLM inference. Skilled in Verilog/SystemVerilog, Synopsys DC, VTR/VPR (C++), and Python performance, energy, and area modeling.

% =====================================================================
% EDUCATION
% =====================================================================
\cvsection{Education}

\entry{Arizona State University}{Tempe, AZ}
\subentry{Ph.D. in Computer Engineering}{08/2024 -- Present}

\entry{Carnegie Mellon University}{Pittsburgh, PA}
\subentry{Master of Science in Electrical and Computer Engineering}{02/2022 -- 12/2024}

\entry{University of Nottingham, Ningbo, China}{Ningbo, China}
\subentry{Bachelor of Engineering in Electrical and Electronics Engineering}{09/2017 -- 07/2021}

% =====================================================================
% INDUSTRY EXPERIENCE
% =====================================================================
\cvsection{Industry Experience}

\entry{Research Associate Intern --- HPE Labs}{Fort Collins, CO}
\subentry{Emerging Accelerator Team}{05/2026 -- Present}
\begin{cvlist}
  \item Developing FPGA--ReRAM in-memory-computing chiplet architectures for LLM inference.
  \item Building system-level performance, energy, and area models for chiplet design-space exploration.
\end{cvlist}

% =====================================================================
% RESEARCH EXPERIENCE
% =====================================================================
\cvsection{Research Experience}

\entry{Graduate Research Assistant --- Arizona State University}{Tempe, AZ}
\subentry{FPGA architecture and CAD, in-memory computing}{08/2024 -- Present}

\project{\emph{NIFA}: Nonlinear IMC Enhanced FPGA for Efficient ML Inference}{International Conference on Field-Programmable Logic and Applications (FPL), 2026}
\begin{cvlist}
  \item Proposed ADC-free ReRAM IMC blocks for FPGAs, with analog CAMs (ACAMs) for nonlinear ops and log-domain attention.
  \item Evaluated the architecture in VTR on ResNet-9, VGG-11, and BERT-Tiny benchmarks (Fmax, wirelength, area).
  \item Ran FPGA-aware DSE of crossbar size and IMC area budget.
  \item Extended a Python IMC-FPGA simulator for end-to-end performance, energy, and area.
  \item Achieved up to 40$\times$/1.9$\times$ energy and 4.1$\times$/2.5$\times$ area efficiency gains over the prior SOTA on CNN and Transformer workloads.
  \item Built a Verilog IP library (GEMM, dynamic MatMul, online softmax, attention), cycle-matched to Python and NumPy models.
\end{cvlist}

\project{Direct BRAM-DSP Paths for Deep Learning on FPGAs}{International Conference on Field-Programmable Technology (FPT), 2026}
\begin{cvlist}
  \item Proposed hardened BRAM-to-DSP direct paths that bypass global routing, at negligible overhead (Synopsys DC, FreePDK45).
  \item Extended the VPR placer (C++) with cross-type placement macros and connectivity-weighted placement priority.
  \item Built a Python place-and-route evaluation flow and evaluated the architecture on FC, Conv, attention, and LSTM benchmarks.
  \item Achieved up to 25\% higher Fmax and 49\% lower wirelength on DL designs.
\end{cvlist}

\project{\emph{RACAM}: Reuse-Aware Bit-Serial DRAM-PIM with Automated Mapping}{International Conference on Parallel Architectures and Compilation Techniques (PACT), 2026}
\begin{cvlist}
  \item Built the initial bit-serial DRAM-PIM GEMM/GEMV performance model on LLMCompass (Python).
  \item Designed the initial tiling and mapping search (array/bank/device tiling, broadcast, loop order).
  \item Implemented bit-serial integer PEs and bank/rank/device broadcast/mux logic in SystemVerilog; synthesized with Synopsys DC.
  \item Built the evaluation pipeline for GPT-3 and Llama-3.1 prefill/decode against GPU baselines.
\end{cvlist}

% =====================================================================
% PUBLICATIONS
% =====================================================================
\cvsection{Publications}

\begin{cvlist}\interlinepenalty=10000  % never split one entry across pages
  \item \textbf{Jiajun Hu}, Ruthwik Reddy Sunketa, Lei Zhao, Archit Gajjar, Luca Buonanno, and Aman Arora. \emph{NIFA: Nonlinear IMC Enhanced FPGA for Efficient ML Inference}. In \emph{International Conference on Field-Programmable Logic and Applications (FPL)}, 2026. \href{https://arxiv.org/abs/2607.15123}{arXiv:2607.15123}
  \item \textbf{Jiajun Hu}, Ruthwik Reddy Sunketa, Andrew Boutros, and Aman Arora. \emph{Boosting FPGA Performance with Direct BRAM-DSP Paths}. In \emph{International Conference on Field-Programmable Technology (FPT)}, 2026 (short paper). \href{https://arxiv.org/abs/2607.05756}{arXiv:2607.05756}
  \item Siyuan Ma, \textbf{Jiajun Hu}, Jeeho Ryoo, Aman Arora, and Lizy Kurian John. \emph{RACAM: Enhancing DRAM with Reuse-Aware Computation and Automated Mapping for ML Inference}. In \emph{International Conference on Parallel Architectures and Compilation Techniques (PACT)}, 2026. \href{https://arxiv.org/abs/2512.09304}{arXiv:2512.09304}
  \item \textbf{Jiajun Hu}, Chetan Choppali Sudarshan, Maxwell Clifford, Vidya Chhabria, and Aman Arora. \emph{CarbonSet: A Dataset to Analyze Trends and Benchmark the Sustainability of CPUs and GPUs}. In \emph{Great Lakes Symposium on VLSI (GLSVLSI)}, 2025, pp.~177--184. \href{https://doi.org/10.1145/3716368.3735235}{doi:10.1145/3716368.3735235}
  \item Chetan Choppali Sudarshan, \textbf{Jiajun Hu}, Aman Arora, and Vidya A. Chhabria. \emph{CarbonPATH: Carbon-aware pathfinding and architecture optimization for chiplet-based AI systems}. arXiv preprint (under review). \href{https://arxiv.org/abs/2603.03878}{arXiv:2603.03878}
\end{cvlist}

\end{document}
